\documentclass[11pt]{article}
\usepackage{fontspec}
\usepackage[french,english]{babel}
\usepackage{amsmath}
\usepackage{amssymb}
\usepackage{amsthm}
\usepackage{mathrsfs}
\usepackage{tikz-cd}
\usepackage{enumitem}

\newcommand{\printedpage}[1]{\marginpar{\footnotesize p.~#1}}
\newcommand{\msp}[1]{\marginpar{\footnotesize [#1]}}
\newcommand{\Fl}{\mathrm{Fl}}
\newcommand{\Is}{\mathrm{Is}}
\newcommand{\Fib}{\mathrm{Fib}}

\begin{document}

\printedpage{54}
% head: A. Grothendieck : Dérivateurs, Ch. II, Édition des Universités de Montpellier II et Paris VII
We are of course interested above all in the closed elements of $\mathcal{P}$, $\mathcal{Q}$, i.e. the subsets $A$, $L$ of $\Fl(\mathcal{M})$, containing the isomorphisms, stable under composition (and $L$ also stable under base change, $A \subseteq W$) such that we have
\[
A = \mathbf{A}_W(L), \qquad L = \mathbf{L}(A) \tag*{(6.8)}
\]
\texttt{[rather }$\mathbf{L}_W(A)$\texttt{]}.

\msp{page 87}

I shall keep the notations $\natural$, $\flat$ of the preceding paragraph. What are the two extreme pairs?

Starting from $A_0 = \Is(\mathcal{M})$ (smallest of the elements of $\mathcal{P}$), we find $\mathbf{L}_W(A_0) = \Fl(\mathcal{M})$ (largest element of $\mathcal{Q}$), and $A_0^{\natural} = \mathbf{A}_W(\Fl(\mathcal{M})) = W_{\mathrm{univ}}$, the set of arrows which are in $W$ and remain so under any base change:
\[
A_0^{\natural} = W_{\mathrm{univ}}, \qquad \beta_0(A_0) \overset{\text{def}}{=} \mathbf{L}_W(A_0) = \Fl(\mathcal{M}), \qquad \text{\scriptsize where $A_0 = \Is(\mathcal{M})$ smallest element of $\mathcal{P}$.} \tag*{(6.9)}
\]
Conversely, let us start from $L_0 = \Is(\mathcal{M})$, the smallest element of $\mathcal{Q}$. We have $\mathbf{A}_W(L_0) = W$ ($W$ being stable under composition and containing the isomorphisms), and $L_0^{\flat} = \mathbf{L}_W(W) = \Fib_W$, the $W$-fibrations:
\[
L_0^{\flat} = \Fib_W, \qquad \alpha_0(L_0) = W, \qquad \text{\scriptsize where $L_0 = \Is(\mathcal{M})$, smallest element of $\mathcal{Q}$.} \tag*{(6.10)}
\]
So the left $\mathcal{R}$-closed subsets of $\Fl(\mathcal{M})$ wander between $W_{\mathrm{univ}}$ and $W$, the right $\mathcal{R}$-closed subsets between $\Fib_W$ and $\Fl(\mathcal{M})$.

\textbf{Proposition 6.1.} \textit{Let $A \in \mathcal{P}$, $L \in \mathcal{Q}$ be such that $A = \alpha_0(L)$, $L = \beta_0(A)$ (i.e. $A = \mathbf{A}_W(L)$, $L = \mathbf{L}_W(A)$). For $L$ to be stable under composition, it is necessary and sufficient that $A$ be stable under base change by the arrows $f \in L$. (More generally, if $L \in \mathcal{Q}$ is any element of $\mathcal{Q}$, and if we assume in} \msp{page 88} \textit{addition that $L$ is stable under composition, then $A = \mathbf{A}_W(L)$ is stable under base change by $L$.)}

Terminology: for $A \in \mathcal{P}$, i.e. $A$ a subset of $\Fl(\mathcal{M})$ contained in $W$ containing the isomorphisms and stable under composition, $\mathbf{L}_W(A)$ is called the \textit{transporter of $A$ into $W$}. (There is a dual notion of \textit{cotransporter.}) It is an element of $\mathcal{Q}$, thus a subset of $\Fl(\mathcal{M})$ containing the isomorphisms and stable under base change. On the other hand, for $L$ in $\mathcal{Q}$, the arrows $f \in \mathbf{A}_W(L)$ are called \textit{$W$-$L$-aspheric,} or the \textit{$W$-$L$-equivalences,} or the \textit{$W$-equivalences that are $L$-universal,} and $\mathbf{A}_W(L)$ is called the \textit{$L$-kernel of $W$}. We could denote it more simply by $W_L$ (when this notation does not conflict with $W_S \subseteq \Fl(\mathcal{M}/S)$ for $S \in \mathrm{Ob}\,\mathcal{M}$). It is a subset of $W$ which (like $W$ itself) contains the isomorphisms and is stable under composition. % [sic] élement -> élément, missing accent

Finally, there is also reason to introduce the \textit{stabilizer} $\mathrm{Stab}(A)$, or transporter of $A$ into $A$, which does not depend on the choice of $W$, but only on the subset $A$ of $\Fl(\mathcal{M})$, containing the isomorphisms and stable under composition. We could also denote it by $\Fib_A$, and call its elements \textit{$A$-fibrations} ($A$ now playing the role of $W$). We have $\mathrm{Stab}(A) \subseteq \mathbf{L}_W(A)$, and it would seem that in % [sic] donné -> donnée, feminine noun
% foot: Version 8 juillet 2001
\printedpage{55}
% head: A. Grothendieck : Dérivateurs, Ch. II, Édition des Universités de Montpellier II et Paris VII
\msp{page 89}
practice, it is these stabilizers, more than the transporters $\mathbf{L}_W(A)$, that are important (when the two do not coincide). Thus in Cat, if $W$ is a basic localizer, the $W$-smooth (resp. $W$-proper) morphisms $\mathrm{Liss}_W$ ($\mathrm{Prop}_W$), which are the stabilizers of the sets $W_{\mathrm{as}}$ ($W_{\mathrm{coas}}$) of the sets of $W$-aspheric (resp. $W$-coaspheric) morphisms, playing an essential role, whereas I do not know how to characterize in a meaningful way the arrows of the transporters $\mathbf{L}_W(W_{\mathrm{as}})$, $\mathbf{L}_W(W_{\mathrm{coas}})$, which contain, in addition to the $W$-smooth (resp. $W$-proper) morphisms the $W$-fibrations $\Fib_W$. I do not know how to characterize a class of morphisms which looks so motley! Likewise, there is reason to introduce and characterize the stabilizer of $W_{\mathrm{bias}}$ (set of $W$-biaspheric arrows), as the ``$W$-neat'' morphisms $\mathrm{Net}_W$, by a description similar to that of the $\mathrm{Liss}_W$, $\mathrm{Prop}_W$, using the bilocalized categories $\backslash\underset{u}{X}/$, $\backslash\underset{v}{Y}/$ , instead of the localized categories $X/x$, $Y/y$ or the colocalized categories $x\backslash X$, $y\backslash Y$. Finally there is reason to determine the stabilizer of the set Equ of Cat-equivalences (which does not depend on the choice of a $W$ - but we have $\mathrm{Equ} \subseteq \mathcal{P}(W)$ for every basic $W$). % [sic] donné -> donnée, feminine noun
% foot: Version 8 juillet 2001

\printedpage{56}
% head: A. Grothendieck : Dérivateurs, Ch. II, Édition des Universités de Montpellier II et Paris VII
\msp{page 90}

\section*{7\quad Cofiltering and locally cofiltering morphisms}

But I do not want to dwell now on these developments, which do not seem urgent to me, and would like to return to the questions connected with the notion of filtering or cofiltering categories.

\textbf{Definition 7.1.} Let $f : X \longrightarrow Y$ be an arrow in Cat. We say that $f$ is \textit{cofiltering} (resp. \textit{locally cofiltering}, resp. \textit{weakly cofiltering} (\footnote[64]{I have changed the terminology: now weakly cofiltering $=$ 0-connected $+$ locally weakly cofiltering, \textit{stable under passage} to $X'$ over $X$, $X'$ étale and $\infty$-connected. Review the terminology in the weakly cofiltering case, which is a treacherous notion, for lack of being stable under passage \texttt{[to]} an $X'$ étale over $X$.}), resp. \textit{locally weakly cofiltering} or again \textit{precofiltering}) if for every $y$ in $Y$, the colocalized category $y\backslash X$ is cofiltering (resp. \dots). (Dually, we say that $f$ is \textit{filtering}, resp. \dots, if the categories $X/y$ are filtering, resp. \dots, i.e. if $f^o : X^o \longrightarrow Y^o$ is cofiltering resp. \dots.) % [sic] Definition -> Définition, missing accent; colocalisé -> colocalisée, agreement with catégorie

\textbf{Remark 7.2.} Notions stable under coétale base change. Contains the isomorphisms and even the Cat-equivalences. The \textit{étale} morphisms are locally filtering.

\textbf{Proposition 7.3.} \textit{Let $f : X \longrightarrow Y$ be a cofiltering morphism (resp. locally cofiltering, resp. weakly cofiltering, resp. precofiltering), and suppose that $Y$ is cofiltering (resp. locally cofiltering, resp. weakly cofiltering, resp. precofiltering). Then $X$ itself is cofiltering (resp. \dots).}

\textsc{Proof}. It is a matter, in the four cases considered, of checking for $X$ the condition $\underbrace{\mathrm{PS}_g1' \;/\; \mathrm{PS}_g1}_{\text{last two cases}} \;/\; \mathrm{PS}_g1' + \mathrm{PS}_g2 \;/\; \mathrm{PS}_g1 + \mathrm{PS}_g2$. But \msp{page 91} for each of the three conditions $\mathrm{PS}_g1$, $\mathrm{PS}_g1'$, $\mathrm{PS}_g2$ separately, we see immediately that if it is satisfied for $Y$ and for the $y\backslash X$, it is satisfied for $X$, whence the proposition.

\textbf{Corollary 7.3.} \textit{The four notions considered in (7.1) are stable under composition.}

\textsc{Proof} immediate, using remark (7.2).
\[
\begin{tikzcd}
X \arrow[d] & z\backslash X \arrow[l] \arrow[d] \\
Y \arrow[d] & z\backslash Y \arrow[l] \arrow[d] \\
Z & z\backslash Z. \arrow[l]
\end{tikzcd}
\]
% foot: Version 8 juillet 2001

\printedpage{57}
% head: A. Grothendieck : Dérivateurs, Ch. II, Édition des Universités de Montpellier II et Paris VII
\textbf{Remark 7.4.}
\begin{enumerate}
\item[1)] Note that the notion of a cofiltering or weakly cofiltering morphism is \textit{not} stable under arbitrary base change, nor even under an open immersion. Example: $e \longrightarrow \boldsymbol{\Delta}^1$, mapping $e$ to 1, is cofiltering,
\[
\begin{tikzcd}
 & e \arrow[d, mapsto] \\
0 \arrow[r] & 1
\end{tikzcd}
\]
but its fiber at 0 is empty, hence is neither cofiltering nor weakly cofiltering. I have not seen an example showing that the notion of a \textit{locally} cofiltering resp. precofiltering arrow is not stable under base change, and would not want to dwell on it, having no need to know the answer.
\item[2)] For $X$ to be cofiltering (resp. \dots), it is necessary and sufficient that $X \longrightarrow e$ be so. Thus the absolute notion is reduced to a relative notion.
\item[3)] If $f : X \longrightarrow Y$ is cofiltering, it is \msp{page 92} $W$-coaspheric (thus $\in W$) for every $W$ satisfying the direct limit axiom Loc (8). In particular, it is $W_0$-aspheric, i.e. a cofiltering functor (and even merely weakly cofiltering) is \textit{cofinal.}
\end{enumerate}
The main reason for my interest in cofiltering arrows is contained in the

\textbf{Proposition 7.4.} \textit{Let $f : X \longrightarrow Y$ in} Cat \textit{be locally cofiltering (resp. locally weakly cofiltering). Then:}
\begin{enumerate}
\item[\textit{a)}] \textit{For every $X'$ étale over $X$, the composite $f' : X' \longrightarrow X \longrightarrow Y$ is locally cofiltering (resp. \dots), and for every factorization of $f'$ as $X' \xrightarrow{g'} Y' \longrightarrow Y$ with $Y' \longrightarrow Y$ étale, $g'$ is locally cofiltering (resp. \dots).}
\item[\textit{b)}] \textit{$\forall\ x$ in $X$, the morphism on the localized categories}
\[
f/x : X/x \longrightarrow Y/x \qquad (y = f(x)) % [sic] Y/x -> Y/y, since y = f(x)
\]
\textit{is cofiltering (resp. weakly cofiltering).}
\item[\textit{c)}] \textit{If $x$, $y$ are as in b), and $U$ is an open subset of $X/x$, denoting by $V$ the open subset of $Y/y$ generated by the image of $f/x$ (i.e. the open subset consisting of the $z$ in $Y/y$ such that $z\backslash(X/x)$ is nonempty), the induced arrow $q : U \longrightarrow V$} \texttt{[is]} \textit{cofiltering (resp. weakly cofiltering), a fortiori cofinal (and even $W$-coaspheric, for every $W$ satisfying the limit axiom Loc (8))} (\footnote[65]{It suffices here to be under the resp'ed hypothesis ($f$ \textit{weakly} locally cofiltering) except for the NB here \texttt{[``the NB here'' is ``(and even ...Loc (8))'']}. % [sic] respée -> intended word unclear (perhaps « respectée » or « resp. »), not a French word
}). % doubt resolved: PDF prints q here; later g is an edition naming inconsistency (recorded in corrections.csv).

\msp{page 93}
\textit{Thus for every} cosheaf \textit{$G$ on $V$, with values in a category $\mathcal{M}$ with $\varinjlim$ , we have} \texttt{[an]} \textit{isomorphism} % [sic] pout -> pour, misprint
\[
\varinjlim_U g^*(G) \xrightarrow{\ \sim\ } \varinjlim_V G,
\]
% foot: Version 8 juillet 2001
\printedpage{58}
% head: A. Grothendieck : Dérivateurs, Ch. II, Édition des Universités de Montpellier II et Paris VII
\textit{and for every presheaf $F$ on $V$, with values in a category $\mathcal{M}$ with $\varprojlim$ , we have} \texttt{[an]} \textit{isomorphism}
\[
\varprojlim_U g^*(F) \xleftarrow{\ \sim\ } \varprojlim_V F.
\]
\end{enumerate}

\textsc{Proof}. I shall first spell out the

\textbf{Lemma 7.5.} \textit{Let $h : X' \longrightarrow X$ be an étale morphism in} Cat. \textit{Then $h$ is locally cofiltering. More precisely, for every $x$ in $X$, we have}
\[
\pi_0(x\backslash X') \xrightarrow{\ \sim\ } \mathrm{Ob}\,X'_x,
\]
\textit{every connected component of $x\backslash X'$ contains one and only one element of the discrete category $X'_x \subseteq x\backslash X'$. If $x'$ is in $X'_x$ (i.e. $x'$ an object of $X'$ such that $h(x') = x$), it is an} initial \textit{object of the connected component which contains it. (Thus this component is cofiltering.)}

This proves (and beyond) the first assertion of 7.4, in view of 7.3 (stability under composition), or because an \msp{page 94} étale category (here $y\backslash X'$) over a locally cofiltering (resp. locally weakly cofiltering) category (here $y\backslash X$) is likewise.
\[
\begin{tikzcd}[row sep=large]
y\backslash X' \arrow[r] \arrow[d, "\text{ét.}"'] \arrow[dr, phantom, "\text{cart.}"] & X' \arrow[d, "\text{ét.}"] \\
y\backslash X \arrow[r] \arrow[d] \arrow[dr, phantom, "\text{cart.}"] & X \arrow[d] \\
y\backslash Y \arrow[r] & Y.
\end{tikzcd}
\]
For the second assertion in a), we can now assume $X' = X$. We must prove that for every $y'$ in $Y'$, $y'\backslash X \overset{\text{def}}{=} X \times_{Y'} y'\backslash Y'$ is locally cofiltering (resp. locally weakly cofiltering). Let $y$ in $Y$ be the image of $y'$, and consider the diagram of three cartesian squares opposite.
\[
\begin{tikzcd}[row sep=large]
y'\backslash X \arrow[r] \arrow[d, "\text{ét.}"] \arrow[dr, phantom, "\textit{cart.}"] & y\backslash X \arrow[r] \arrow[d, "\text{ét.}"] \arrow[dr, phantom, "\textit{cart.}"] & X \arrow[d, "\text{ét.}"] \\
y'\backslash Y' \arrow[r] & y\backslash Y' \arrow[r] \arrow[d] \arrow[dr, phantom, "\textit{cart.}"] & Y' \arrow[d] \\
 & y\backslash Y \arrow[r] & Y
\end{tikzcd}
\]
Lemma 7.5 implies that $y'\backslash Y' \longrightarrow y\backslash Y'$ is the inclusion into $y\backslash Y'$ of a connected component of $y\backslash Y'$, so it is an open immersion. So $y'\backslash X \longrightarrow y\backslash X$ is an immersion
% foot: Version 8 juillet 2001

\end{document}
